\newcommand{\docshorttitle}{SUPPL.\ MEMO P\&A --- ``CUSTOMER'' DOCTRINE \S 425.17(c) (CGC-25-631802)}
\newcommand{\doctitleblock}{PLAINTIFF'S SUPPLEMENTAL MEMORANDUM OF POINTS AND AUTHORITIES IN FURTHER OPPOSITION TO DEFENDANT'S SPECIAL MOTION TO STRIKE --- ``CUSTOMER'' DOCTRINE UNDER \S 425.17(c) AND THE THIRD-PARTY CLAIMANT}
\newcommand{\docsubblock}{(Cal.\ Rules of Court, rules 3.1113, 3.1306(c); supplemental to Plaintiff's previously-filed Opposition; cross-noticed for hearings of May 13, 2026 and May 19, 2026, Dept.\ 302)}
% Pleading-paper preamble for APRIL-24-2026 cover sheets and oral-argument
% outlines (CGC-25-631801 and CGC-25-631802). Matches house style from
% APRIL-23-SAC-DEFENSE/court-pdfs/REPLY-SAC-PREAMBLE.tex, simplified for
% single-to-three-page artifacts that do not require TOC/TOA infrastructure.
%
% Cross-hearing caption (CAPTION-801/802.tex): use \CaptionDocTitleBlock,
% \CaptionDocSubBlock, and the crosshearingsschedule environment (defined below)
% so long \doctitleblock text and hearing dates stay inside the 3in minipage.
\documentclass{article}
\usepackage[utf8]{inputenc}
\usepackage{graphicx}
\usepackage{geometry}
    \geometry{top=1in, bottom=2.2in, left=1.0in, right=1.0in, textheight=336pt}
    \setlength{\footskip}{70pt}
\usepackage{ulem}
\usepackage{fancyhdr}
    \renewcommand{\headrulewidth}{0pt}
    \renewcommand{\footrulewidth}{0.4pt}
\usepackage{parskip}
    \setlength{\parskip}{0mm}
\usepackage{quoting}
    \quotingsetup{leftmargin=0.5in, rightmargin=0.5in, vskip=-1.5mm}
    \makeatletter
        \g@addto@macro\quoting\singlespacing
        \g@addto@macro\quoting{\vspace{-2mm}}
    \makeatother
\usepackage{mathptmx}
\usepackage{amssymb}
\renewcommand{\rmdefault}{ptm}
\renewcommand{\normalsize}{\fontsize{12}{12}\selectfont}
\newcommand*{\ptm}{\fontfamily{ptm} \selectfont \fontsize{12}{12}\selectfont}
\usepackage{setspace}
    \doublespacing
\newcommand{\tab}{\hspace*{.0in}}
\newenvironment{tightcenter}{%
    \setlength\topsep{0pt}
    \setlength\parskip{0pt}
    \begin{center}
}{%
    \end{center}
}
\usepackage{tikz}
\usetikzlibrary{calc}
% Pleading-paper vertical double rule at the left margin.
\AddToHook{shipout/background}{%
    \begin{tikzpicture}[remember picture, overlay]
        \draw[line width=0.5pt] ($(current page.north west) + (0.75in,0)$) -- ($(current page.south west) + (0.75in,0)$);
        \draw[line width=0.5pt] ($(current page.north west) + (0.80in,0)$) -- ($(current page.south west) + (0.80in,0)$);
    \end{tikzpicture}%
}
\usepackage[left, pagewise]{lineno}
\setlength{\linenumbersep}{0.35in}
\renewcommand{\linenumberfont}{\normalfont\small}
\usepackage{titlesec}
\usepackage[hidelinks]{hyperref}
\usepackage{bookmark}
\usepackage{textcomp}
\usepackage{longtable}
\usepackage{booktabs}
\usepackage{array}
\usepackage{multirow}
\usepackage{calc}
% --- Cross-hearing caption cover (3in minipage): compact title + 3-column hearing table ---
% Use on first page only; keeps long \doctitleblock / \docsubblock on one sheet with party caption.
\newcommand{\CaptionDocTitleBlock}{%
  \begingroup
  \fontsize{10}{11.5}\selectfont
  \bfseries
  \raggedright
  \setlength{\parskip}{0pt}%
  \noindent\doctitleblock\par
  \endgroup
}
\newcommand{\CaptionDocSubBlock}{%
  \begingroup
  \footnotesize
  \raggedright
  \setlength{\parskip}{0pt}%
  \noindent\docsubblock\par
  \endgroup
}
% #1 = department number only (e.g.\ 301 or 302) for the line ``Dept.\ #1''.
\newenvironment{crosshearingsschedule}[1]{%
  \par\begingroup
  \footnotesize
  \setlength{\tabcolsep}{2.2pt}%
  \renewcommand{\arraystretch}{1.04}%
  \setlength{\parskip}{0pt}%
  \setlength{\parindent}{0pt}%
  \noindent\textbf{Cross-noticed for pending defense hearings (this case, Dept.\ #1):}\par\vspace{2pt}%
  \begin{tabular}{@{}>{\raggedright\arraybackslash}p{0.11\linewidth}>{\raggedright\arraybackslash}p{0.34\linewidth}>{\raggedright\arraybackslash}p{0.53\linewidth}@{}}%
  \textbf{\#} & \textbf{Motion} & \textbf{Date / dept / judge} \\ \hline
}{%
  \end{tabular}\par\endgroup\smallskip
}
% Pandoc pipe-table column widths use \real{#}.
\newcommand{\real}[1]{\pgfmathparse{#1}\pgfmathresult}
\titleformat{\section}{\fontsize{13}{13}\selectfont\normalfont\bfseries\uppercase}{}{0em}{}
\titleformat{\subsection}{\fontsize{12}{12}\selectfont\normalfont\bfseries}{}{0em}{}
\titleformat{\subsubsection}{\normalfont\bfseries}{}{0em}{}
\titlespacing*{\section}{0pt}{6pt}{2pt}
\titlespacing*{\subsection}{0pt}{4pt}{2pt}
\titlespacing*{\subsubsection}{0pt}{2pt}{2pt}
\newcommand{\calrules}[1]{Cal. Rules of Court, rule #1}
\newcommand{\ccp}[1]{Code Civ. Proc., \textsection~#1}
\newcommand{\evidcode}[1]{Evid. Code, \textsection~#1}
\providecommand{\tightlist}{%
  \setlength{\itemsep}{0pt}\setlength{\parskip}{0pt}}
% Footer: page N of total + document short-title.
\usepackage{lastpage}
\pagestyle{fancy}
\fancyhf{}
\fancyfoot[C]{\thepage{} of \pageref{LastPage} \\[1pt] \footnotesize \docshorttitle}
\fancyfoot[R]{}

\begin{document}
% Cross-hearing caption for CGC-25-631802 (Dept. 302, Hon. Quinn).
\nolinenumbers
\begin{singlespace}
\noindent Franciscus Dylan Rosario \\
7624 Melody Drive \\
Rohnert Park, CA 94928 \\
Tel: 650.488.7510 \\
E: soltrinox@gmail.com \\
Plaintiff, In Pro Per
\end{singlespace}
\vspace*{9mm}
\begin{tightcenter}
\textbf{SUPERIOR COURT OF THE STATE OF CALIFORNIA} \\
\textbf{COUNTY OF SAN FRANCISCO}
\end{tightcenter}
\vspace*{3mm}
\begin{minipage}[t]{3in}
    \raggedright
    \textbf{FRANCISCUS DYLAN ROSARIO,} \\ \textbf{PLAINTIFF}, \\
    \hspace{1cm} v. \\
    \small
    CSAA Insurance Exchange; and Does 1 through 50, inclusive;\\
    \normalsize
    \textbf{DEFENDANTS}
    \makebox[3in]{\hrulefill}
\end{minipage}%
\hspace{3mm}%
\begin{minipage}[t]{0.5pt}
    \vspace{0pt}
    \rule{0.5pt}{2.42in}
\end{minipage}%
\hspace{3mm}%
\begin{minipage}[t]{3in}
    \raggedright
    \noindent\textbf{Case No.: CGC-25-631802}\par\smallskip
    \CaptionDocTitleBlock\smallskip
    \CaptionDocSubBlock\smallskip
    \begin{singlespace}
    \setlength{\parindent}{0pt}
    \begin{crosshearingsschedule}{302}
    (1) & MFL / leave (SAC) & \textbf{Apr.\ 30, 2026} \\
    (2) & Anti-SLAPP (\ccp{425.16}) & \textbf{May 13, 2026, 9:30 a.m.};\ Dept.\ 302;\ Hon.\ Joseph M.\ Quinn \\
    (3) & Strike (\ccp{436}) & \textbf{May 19, 2026};\ Dept.\ 302;\ Hon.\ Joseph M.\ Quinn \\
    (4) & Demurrer & Hrng.\ date \textbf{[TBD]} (reserved) \\
    \end{crosshearingsschedule}
    \noindent Complaint filed: Dec.\ 23, 2025 \\
    \noindent FAC filed: Feb.\ 25, 2026 \\
    \end{singlespace}
\end{minipage}
\vspace*{8mm}
\newpage
\linenumbers


\section*{I. INTRODUCTION AND VEHICLE}

This Supplemental Memorandum of Points and Authorities is respectfully submitted in further opposition to Defendant CSAA Insurance Exchange's special motion to strike under Code of Civil Procedure section 425.16. It is offered as a supplement to --- and not as a substitute for --- Plaintiff's previously-filed Opposition and the companion Supplemental Memorandum on the structural sequencing of \ccp{425.17(c)} filed concurrently in this case. The vehicle is \calrules{3.1113} (memorandum of points and authorities) read together with \calrules{3.1306(c)} (lodging of supplemental authority for hearings).

This memorandum is directed at one discrete question that the defense's papers attempt to manufacture: whether Plaintiff, as a third-party claimant, is a ``customer'' of CSAA's regulated insurance-claims-handling product within the meaning of \ccp{425.17(c)(2)}. The defense contends that, because Plaintiff is not a named insured under a CSAA policy, Plaintiff is not within the audience the carve-out protects, and therefore Defendant's special motion to strike may go forward.

That argument collides with three independent bodies of California law: (i) the Insurance Commissioner's formal regulations at 10 California Code of Regulations section 2695.1 et seq., which expressly define ``claimant'' to include both first party and third party claimants and impose the regulatory floor on ``all persons'' and ``all claims''; (ii) Insurance Code section 790.03(h), which uses the term ``claimants'' --- without qualification --- in twelve of sixteen subdivisions and uses the disjunctive ``insureds or claimants'' to confirm that both categories are protected; and (iii) the disjunctive structure of \ccp{425.17(c)(1)} and (c)(2), as construed by \textit{Simpson Strong-Tie Co.\ v.\ Gore} (2010) 49 Cal.4th 12. Each is independently dispositive. Together they leave no room for the defense's ``no-customer'' theory.

\section*{II. ISSUE PRESENTED}

Whether a third-party claimant whose claim CSAA handled, investigated (or represented that it investigated), and ultimately denied --- on CSAA's own claim form and under its own claim number --- is a ``customer'' of CSAA's regulated insurance-claims-handling product within the meaning of Code of Civil Procedure section 425.17(c)(2).

\textbf{Plaintiff's answer:} Yes. The Insurance Commissioner has, by formal rulemaking, classified third-party claimants as members of the regulated class; the Legislature has expressly recognized claimants generally as a protected category in section 790.03(h); and \ccp{425.17(c)(2)} is, in any event, disjunctive, so that the audience element is independently satisfied by branches (b) and (c) of the subdivision regardless of how the term ``customer'' in branch (a) is construed.

\section*{III. ARGUMENT}

\subsection*{A. The Insurance Commissioner's Operative Definitions Under Title 10 Treat First-Party and Third-Party Claimants as a Single Regulated Class.}

The Insurance Commissioner has, by formal rulemaking, defined the operative term ``claimant'' for purposes of California claims-handling regulation. The definition is unambiguous:

\begin{quoting}
``\,`Claimant' means a first party claimant, a third party claimant, or both and includes such claimant's designated legal representative and includes a member of the claimant's immediate family designated by the claimant.''
\end{quoting}
\begin{flushright}
\small (Cal.\ Code Regs., tit.\ 10, \S~2695.2(c).)
\end{flushright}

\noindent The companion definitions resolve any residual doubt:

\begin{quoting}
``\,`First party claimant' means any person asserting a right under an insurance policy as a named insured, other insured, or beneficiary under the terms of that insurance policy.''
\end{quoting}
\begin{flushright}
\small (Cal.\ Code Regs., tit.\ 10, \S~2695.2(s).)
\end{flushright}

\begin{quoting}
``\,`Third party claimant' means any person asserting a claim against any person under an insurance policy.''
\end{quoting}
\begin{flushright}
\small (Cal.\ Code Regs., tit.\ 10, \S~2695.2(t).)
\end{flushright}

\noindent The scope and applicability provisions are correspondingly broad:

\begin{quoting}
``The regulations which follow set forth the minimum standards for the settlement of claims and shall be applicable to the handling or settlement of \textbf{all claims}, except as specifically excluded.''
\end{quoting}
\begin{flushright}
\small (Cal.\ Code Regs., tit.\ 10, \S~2695.1(a) (emphasis added).)
\end{flushright}

\begin{quoting}
``These regulations are applicable to \textbf{all persons} and to \textbf{all insurance policies}, insurance contracts, and certificates, except as specifically excluded.''
\end{quoting}
\begin{flushright}
\small (Cal.\ Code Regs., tit.\ 10, \S~2695.1(d) (emphasis added).)
\end{flushright}

The Commissioner did not carve a third-party exemption out of the regulatory floor; he extended the floor to ``all persons'' and to ``all claims.'' The defense's ``no-privity'' theory cannot be reconciled with this text. The Commissioner has, on formal authority delegated by the Legislature under Insurance Code section 790.10, classified third-party claimants as members of the regulated class.

The substantive duties imposed on the regulated class by Title 10 confirm the point. Section 2695.7(b) of Title 10 imposes the duty to accept or deny ``the claim'' within forty calendar days of receiving proof of claim, and section 2695.7(d) imposes the duty to ``conduct and diligently pursue a thorough, fair and objective investigation'' of every claim. Neither provision distinguishes between first-party and third-party postures. Both run to ``the claimant'' as defined in section 2695.2(c) --- that is, to first-party and third-party claimants as a unitary category.

\subsection*{B. Insurance Code Section 790.03(h) Uses the Term ``Claimants'' Without Qualification and Treats Insureds and Claimants as Coextensive Protected Categories.}

The regulations summarized above implement Insurance Code section 790.03(h), the operative text of which is reproduced verbatim in the Supplemental Request for Judicial Notice filed concurrently in this case. For purposes of the customer-doctrine question, what matters is the Legislature's own usage of the word ``claimants'' across the sixteen enumerated unfair claims settlement practices.

Twelve of the sixteen subdivisions reach communications and conduct directed at claimants generally. Specifically:

\begin{itemize}
\tightlist
\item \textbf{\S~790.03(h)(1)}: ``Misrepresenting to \textbf{claimants} pertinent facts or insurance policy provisions relating to any coverages at issue.''
\item \textbf{\S~790.03(h)(2)}: ``Failing to acknowledge and act reasonably promptly upon \textbf{communications with respect to claims} arising under insurance policies.''
\item \textbf{\S~790.03(h)(3)}: ``Failing to adopt and implement reasonable standards for the \textbf{prompt investigation and processing of claims} arising under insurance policies.''
\item \textbf{\S~790.03(h)(5)}: ``Not attempting in good faith to effectuate prompt, fair, and equitable settlements of \textbf{claims in which liability has become reasonably clear}.''
\item \textbf{\S~790.03(h)(10)}: ``Making known to \textbf{insureds or claimants} a practice of the insurer of appealing from arbitration awards in favor of \textbf{insureds or claimants} \dots''
\item \textbf{\S~790.03(h)(11)}: ``Delaying the investigation or payment of claims by requiring \textbf{an insured, claimant, or the physician of either} \dots''
\item \textbf{\S~790.03(h)(13)}: ``\dots~for the \textbf{denial of a claim} or for the offer of a compromise settlement.''
\item \textbf{\S~790.03(h)(14)}: ``Directly advising \textbf{a claimant} not to obtain the services of an attorney.''
\item \textbf{\S~790.03(h)(15)}: ``Misleading \textbf{a claimant} as to the applicable statute of limitations.''
\end{itemize}

Two structural observations follow. \textbf{First}, the Legislature uses the term ``claimants'' repeatedly and without qualification. There is no ``first-party-only'' subdivision among (h)(1), (h)(2), (h)(3), (h)(5), (h)(13), (h)(14), or (h)(15); the duty runs to ``claimants'' categorically. \textbf{Second}, where the Legislature wished to distinguish between insureds and claimants, it did so expressly --- and in the disjunctive. The phrase ``insureds or claimants'' (subdivision (h)(10)) and the phrase ``an insured, claimant, or the physician of either'' (subdivision (h)(11)) make plain that the Legislature understood ``insureds'' and ``claimants'' as related but distinct protected categories. The defense's reading would write that protection out of the statute. If the Legislature had wanted to limit the unfair-practices floor to insureds, it would have written ``insureds'' wherever it wrote ``claimants.'' It did not.

The conclusion is unavoidable: section 790.03(h) treats third-party claimants as a protected category --- as customers of the regulated claims-handling product --- coextensive with insureds for purposes of the duties of communication, investigation, and prompt, fair, and equitable settlement.

\subsection*{C. The California Supreme Court's Trilogy of \textit{Royal Globe} / \textit{Moradi-Shalal} / \textit{Zhang} Confirms That the Regulatory Duty Toward Third-Party Claimants Survived the Contraction of the Private Remedy.}

The trilogy of California Supreme Court decisions construing section 790.03(h) and its interaction with private remedies is consistent on the substantive duty. \textit{Royal Globe Insurance Co.\ v.\ Superior Court} (1979) 23 Cal.3d 880 held that a third-party claimant could maintain a private action for violation of section 790.03(h), explicitly grounding the holding in the recognition that ``claimants'' are a category of persons whom the Legislature intended to protect.

\textit{Moradi-Shalal v.\ Fireman's Fund Ins.\ Companies} (1988) 46 Cal.3d 287 overruled \textit{Royal Globe} on the question of the implied private right of action under section 790.03(h). But \textit{Moradi-Shalal} was explicit that it was overruling \textit{Royal Globe} only on the remedy question, not on the underlying duty. The Court emphasized that section 790.03(h) ``is intended to regulate insurance industry settlement practices,'' that the Insurance Commissioner retains administrative authority to enforce it, and that carriers remain bound by the duties the statute and its implementing regulations impose. (\textit{Moradi-Shalal}, 46 Cal.3d at pp.\ 304--305.)

\textit{Zhang v.\ Superior Court} (2013) 57 Cal.4th 364 then confirmed that \textit{Moradi-Shalal} did not bar a third-party claimant from asserting independent statutory claims (there, under Bus.\ \& Prof.\ Code \S~17200) premised on conduct that violates the standards of section 790.03(h). The Court reasoned that \textit{Moradi-Shalal} addressed only the question whether section 790.03(h) itself created an implied private right of action, not whether the conduct the statute prohibits could be the predicate for an unfair-competition claim. (\textit{Zhang}, 57 Cal.4th at pp.\ 380--384.)

The lesson of this trilogy for the present case is straightforward: even after \textit{Moradi-Shalal}, the regulatory floor that defines third-party claimants as members of a single protected class with insureds remains the operative law of California. The defense cannot claim that \textit{Moradi-Shalal} eliminated third-party claimants from the protective scope of the statute. \textit{Moradi-Shalal} eliminated only the implied private right of action under section 790.03(h) itself; it did not eliminate the duty.

The continuing operation of the duty is also confirmed by \textit{Egan v.\ Mutual of Omaha Ins.\ Co.} (1979) 24 Cal.3d 809, 818--819 (recognizing the insurer's duty to ``investigate the claim thoroughly'' as a baseline expectation of the insurance product) and by \textit{Wilson v.\ 21st Century Ins.\ Co.} (2007) 42 Cal.4th 713, 720--721 (reaffirming that an insurer cannot defend by asserting a ``genuine dispute'' while having failed to conduct the thorough and fair investigation that the statute and regulations require).

\subsection*{D. Code of Civil Procedure Section 425.17(c)'s Use of ``Customer'' Is Functional, Not Contractual; the Audience Prong Is Triply Disjunctive.}

The text of \ccp{425.17(c)} forecloses the defense's narrow reading of ``customer.''

\begin{quoting}
``(2) The intended audience is \textbf{an actual or potential buyer or customer}, or \textbf{a person likely to repeat the statement to, or otherwise influence}, an actual or potential buyer or customer, or \textbf{the statement or conduct arose out of or within the context of a regulatory approval process, proceeding, or investigation} \dots''
\end{quoting}
\begin{flushright}
\small (\ccp{425.17(c)(2)} (emphasis added).)
\end{flushright}

The audience element is satisfied where any one of three independent branches is met: (a) actual or potential buyer or customer; (b) person likely to repeat or influence such a customer; or (c) conduct arising within the context of a regulatory proceeding or investigation. Each branch is independently sufficient.

Even if the defense's narrow reading of branch (a) were somehow correct (it is not), branches (b) and (c) independently apply on the face of this record. A third-party claimant is unquestionably ``a person likely to \dots~influence'' the resolution of the claim. And every claim communication issued by a California-licensed insurer is conduct that ``arose out of or within the context of a regulatory \dots~investigation'' --- namely, the claims-handling investigation that the Commissioner's regulations require under sections 2695.5 and 2695.7. The defense's brief does not engage with branches (b) and (c); but those branches are independently sufficient to satisfy the audience element.

The California Supreme Court's construction of \ccp{425.17(c)} reinforces the functional, not contractual, character of the term ``customer.'' \textit{Simpson Strong-Tie Co.\ v.\ Gore} (2010) 49 Cal.4th 12, 28--32, instructs that section 425.17(c) is to be read consistently with its remedial purpose: to prevent commercial actors from using the anti-SLAPP statute to extract commercial-conduct litigation from the ordinary civil docket. \textit{JAMS, Inc.\ v.\ Superior Court} (2016) 1 Cal.App.5th 984 and \textit{Stewart v.\ Rolling Stone LLC} (2010) 181 Cal.App.4th 664 reinforce that the inquiry under section 425.17(c) is functional, focused on whether the defendant's conduct was part of the commercial channel by which it delivered its product, and not on the formal contractual labels parties might attach to the transaction.

\subsection*{E. The CSAA Denial Letter Itself Establishes Customer Status: Insurers Do Not Deny Claims to Non-Customers.}

The clearest factual demonstration that the defense's ``no-customer'' theory is incompatible with its own conduct is the denial communication CSAA issued to Plaintiff. By the act of issuing a denial --- on its own claim form, under its own claim number, signed by its own claims personnel --- CSAA acknowledged precisely the duty whose existence its anti-SLAPP papers now seek to disclaim.

A California-licensed insurer does not issue claim denials to non-customers. The very act of issuing a denial communicates, on the face of the letter: (i) that the recipient has presented a claim under or against an insurance policy administered by the insurer; (ii) that the insurer has, or represents that it has, conducted an investigation of that claim under its regulatory obligations; (iii) that the insurer has reached a determination disposing of the claim; and (iv) that the insurer is communicating the disposition to the claimant in fulfillment of its duty under Cal.\ Code Regs., tit.\ 10, sections 2695.7(b) and 2695.7(d). Each of those propositions is, by the insurer's own conduct, an acknowledgment that the recipient is within the regulated class --- that is, a customer of the insurance-claims-handling product.

The defense cannot have it both ways. It cannot, on the one hand, issue a denial letter that depends for its meaning on the recipient being within the regulated class, and on the other hand, when sued, deny that the recipient was ever within that class to begin with. That posture is the precise sort of post-hoc re-characterization that section 425.17(c) was enacted to forbid.

\subsection*{F. The Privilege of California Licensure: A Statutorily Created Duty Cannot Be Disclaimed by Re-Characterization.}

California's regulation of the insurance industry rests on the proposition that the privilege of writing insurance in this State is conditional. Insurance Code section 700 et seq.\ requires that no person may transact insurance business in California without a certificate of authority from the Commissioner. The Commissioner's continued grant of that certificate is conditioned on the insurer's compliance with the duties the Legislature has imposed --- including the duties of fair claims handling enumerated in section 790.03(h) and operationalized in Title 10.

In substance: the right to operate as an insurance exchange in California is conferred in exchange for the obligation to handle every claim, first-party and third-party alike, in accordance with the regulatory floor. A carrier that wishes to disclaim third-party claims-handling duties is, in substance, asking the State for a unilateral modification of the bargain on which its certificate of authority was issued. That request is foreclosed by the statutes and regulations themselves.

The point bears directly on the section 425.17(c) analysis. The carve-out exists, in substantial part, to prevent a regulated commercial actor from leveraging the procedural mechanics of section 425.16 to escape the regulatory consequences of the very conduct that defines its regulated business. To allow CSAA to characterize a regulatorily-required denial communication as ``free speech'' --- and thereby to extract the consequent litigation from the ordinary civil docket --- is to permit precisely the abuse that section 425.17(a) condemns.

\section*{IV. APPLICATION TO CGC-25-631802}

The sole defendant is CSAA --- a California-licensed insurer. Plaintiff, on the face of the operative pleadings and on the face of CSAA's own denial communication, is a third-party claimant who presented a claim under or against a CSAA-administered policy. The Title 10 regulations --- sections 2695.2(c), 2695.2(t), 2695.4(a), 2695.5, 2695.7(a), 2695.7(b), and 2695.7(d) --- placed CSAA under the duties of acknowledgment, investigation, communication, and prompt, fair, and equitable disposition with respect to Plaintiff's claim. CSAA's pre-litigation conduct --- including the issuance of a denial communication on its own claim form --- was conduct made in the course of delivering the regulated insurance-claims-handling product to Plaintiff. The audience for that conduct was Plaintiff himself, a person within the regulated class as defined by the Commissioner's own definitions. Section 425.17(c) therefore applies, both branches of subdivision (c)(1) are satisfied, and all three branches of subdivision (c)(2) are independently satisfied. The special motion to strike must be denied at the gateway. (See \textit{Simpson Strong-Tie}, 49 Cal.4th at pp.\ 28--32.)

\section*{V. CONCLUSION}

The defense's effort to evade Code of Civil Procedure section 425.17(c) by re-characterizing Plaintiff as a non-customer of CSAA's regulated insurance-claims-handling product cannot be reconciled with the operative definitions in 10 California Code of Regulations sections 2695.2(c), 2695.2(s), 2695.2(t), 2695.1(a), and 2695.1(d); the substantive duties in sections 2695.4(a), 2695.5, 2695.7(a), 2695.7(b), and 2695.7(d); the statutory text of Insurance Code section 790.03(h); the trilogy of \textit{Royal Globe} / \textit{Moradi-Shalal} / \textit{Zhang}; the substantive content of the carrier's duty to investigate as recognized in \textit{Egan} and \textit{Wilson}; the disjunctive structure of \ccp{425.17(c)(1)} and (c)(2) as construed in \textit{Simpson Strong-Tie}, \textit{JAMS}, and \textit{Stewart}; or CSAA's own conduct in issuing a denial letter on its own claim form.

California has not carved out a third-party exemption from the duties of investigation and fair settlement that condition the privilege of writing insurance in this State. Plaintiff respectfully requests that the Court deny the special motion to strike at the gateway under \ccp{425.17(c)}, without reaching the merits-stage probability-of-prevailing inquiry, with the consequence that no interlocutory appeal lies under \ccp{425.17(e)} and the ordinary civil docket may proceed.

\input{SIGBLOCK.tex}
\end{document}
